\subsection{Architectures}
When designing a FFNN, we need to make several choices regarding the architecture of the network, such as:
\begin{itemize}
    \item \textbf{Number of Layers}: how deep should the network be?
    \item \textbf{Number of Neurons per Layer}: how many neurons should each layer have?
    \item \textbf{Expressiveness of the Network}: how complex should the network be to capture the underlying structure of the data?
\end{itemize}
We know that a 2-layer Neural network are universal function approximators, which means that they can approximate any continuous 
function to an arbitrary degree of accuracy, given enough neurons in the hidden layer.
This implies that we could use a shallow network with a single hidden layer, by increasing the number of neurons in that layer (wide network).

However, in practice, it's not guaranteed that our training algorithm will find the optimal parameters that allow the network to approximate the target function, 
as we have to take into consideration the problems in optimization and generalization.
It has been shown that deeper networks can generally capture more complex functions and achieve better generalization performance than shallower networks.

However, increasing the depth of the network is not always the best solution.
Other approaches instead focused on architecutral design choices rather than depth. In particular, this idea was first explored in the context of Convolutional
Neural Networks (CNNs), where the use of convolutional layers allows to get better performance with fewer parameters compared to FFNNs with fully connected layers.